Hlavním cílem této diplomové práce bylo navrhnout jednoduchého autonomního agenta, který se dokáže za pomoci strojového učení a změny vybavení adaptovat v dynamickém prostředí a snadněji porazit protihráče.

Nejprve jsem popsal existující AI agenty ve hrách a nejčastěji používané techniky a algoritmy. Mezi zmíněnými agenty byli AlphaStar, OpenAI Five a Voyager; z algoritmů a technik pak PPO, Deep Q-Network, Self-Play a Population-Based Training.

V další části jsem uvedl technologie použité při tvorbě samotné hry a agenta — např. Unreal Engine či plugin NevarokML. Popsán byl také návrh hry i samotného agenta.

Závěrečná část práce se věnovala implementaci jednotlivých modulů hry, jako jsou hráčská postava, zbraně, systém kol, získávání vybavení (zbraní a brnění) a především implementaci, trénování a aplikaci agenta. Tento agent využívá logiku vytvořenou pomocí strojového učení a je schopen pohybu, rotace, útoku a změny výbavy.

První verze agenta zvládala pouze základní pohyb a jejím hlavním účelem bylo osvojení si práce s pluginem NevarokML. Druhá verze, která představovala první plně funkční implementaci, vykazovala částečnou úspěšnost — agent se občas zcela zasekl a nevěděl, jak reagovat, nebo se místo udržování odstupu příliš přibližoval k nepříteli. Finální verze agenta, navzdory několika drobným nedostatkům, nejlépe prokázala schopnost adaptace v dynamickém prostředí a na chování hráče. Dokáže se pohybovat, otáčet, útočit a po každém kole měnit svou výbavu.